\SourceLeaf{131}
Wo sieben oder zehn Dezimalstellen ausreichen, sind die gewöhnlichen oder die \textsc{Vlacq}schen, auch von \textsc{Vega} herausgegebenen Tafeln freilich bei Weitem bequemer als die \textsc{Leonelli}schen. Tatsächlich sind jene ausnahmslos selbst für die genauesten astronomischen oder sonstigen auf die physische Welt bezogenen Rechnungen genauer als nötig. Bei analytischen Rechnungen oder bei der erstmaligen Erstellung von Tafeln wünscht man jedoch öfter eine größere Genauigkeit. Sofern man nicht mehr als 13 oder 18 Ziffern verlangt, sind \textsc{Leonellis} Tafeln dafür unstreitig das Brauchbarste und Bequemste.

\textsc{Delambres} Bericht zufolge hat \textsc{Brigg} genau dieselbe Methode in seiner \emph{arithmetica logarithmica} dargestellt und eine ähnliche Hilfstafel auf 14 Dezimalstellen berechnet. Die französischen Geometer selbst haben jedoch nur ein Exemplar dieses seltenen Werkes auftreiben können, in dem gerade dieser Teil fehlte. Das Verdienst der eigenen Erfindung und der eigenen Berechnung bleibt \textsc{Leonelli} also ungeschmälert, und es ist gerecht, ihm besonders für die Berechnung den gebührenden Dank auszusprechen.

Den Zweck der zweiten Abhandlung haben wir oben bereits angegeben. Nun wollen wir sehen, wie \textsc{Leonelli} ihn erreichen will. Er schlägt eine Tafel mit drei Spalten vor. Der Kürze halber wollen wir die zusammengehörigen Einträge dieser drei Spalten mit $P$, $Q$, $R$ bezeichnen und auf \textsc{Leonellis} umständliche und unnötige Terminologie verzichten. Diese Größen sollen so voneinander abhängen, dass, wenn man die Zahl mit dem Logarithmus $P$ mit $x$ bezeichnet,
\[
Q = \log\left(1+\frac{1}{x}\right), \qquad R = \log(1+x)
\]
gilt. Daher ist stets $R = P+Q$. Ferner soll diese Tafel nicht nach $x$, sondern nach $P$ geordnet sein, also soll $P$ gleichmäßig von null bis ins Unendliche wachsen, oder vielmehr so lange, bis $Q$ als verschwindend angesehen werden kann. Der Gebrauch einer solchen Tafel ist leicht zu überblicken. Soll aus $\log a$ und $\log b$ der Logarithmus von $a+b$ bestimmt werden, so geht man unter der Voraussetzung, dass $a$ größer als $b$ ist, mit $\log a - \log b$ in die erste Spalte ein oder setzt diese Differenz $= P$. Dann gilt offenbar
\[
\log(a+b) = Q + \log a = R + \log b.
\]

